% TOP_PROBLEM: {"id": 2308005, "title": "Research Problems in Function Theory — Problem 8.5", "category_id": 8, "category": {"id": 8, "name": "analysis", "display_name": "Analysis", "description": "Limits, continuity, calculus, and function theory.", "slug": "analysis", "order_index": 8, "created_at": "2026-07-31T15:26:25.671Z"}, "status": "open"}
% FIRST_POSTED: null
% ENTRY_KIND: statement_only
% Imported from: batch06.tex; UnsolvedMath ID 2308005
% Compile twice: lualatex 2308005.tex
\documentclass[11pt]{article}
\usepackage{fontspec}
\setmainfont{Latin Modern Roman}
\usepackage{amsmath,amssymb,mathtools,unicode-math}
\setmathfont{Latin Modern Math}
\usepackage{xurl,hyperref}
\hypersetup{hidelinks,pdftitle={UnsolvedMath Problem 2308005}}
\newcommand{\sref}[2]{\hyperlink{source:#1:#2}{[S#2]}}
\newcommand{\cataloguescope}[1]{\paragraph{Mathematical question.}#1}
\begin{document}
% UnsolvedMath ID 2308005; source catalogue code AMR-022-8005
\setcounter{section}{2308004}
\section{Research Problems in Function Theory — Problem 8.5}\label{2308005}
{\small\sffamily UnsolvedMath ID 2308005 · Analysis}
\subsection{Definitions and mathematical statement}

\paragraph{Shared notation.}$\mathbb N=\{1,2,\ldots\}$, and $\mathbb Z,\mathbb Q,\mathbb R,\mathbb C$ denote the integers, rationals, reals and complex numbers. Any other conventions are specified locally.


For a connected open $G\subset\mathbb C$, let $H(G)=\mathcal O(G)$ with pointwise addition and multiplication, regarded as a structure in the ring language $\{0,1,+,-,\cdot\}$. Two such structures are elementarily equivalent if they satisfy the same first-order sentences in that language. This differs from $H_{\mathbb C}(G)$ as a complex algebra, whose language additionally names the complex scalar constants. Conformal or anticonformal equivalence means the existence of a holomorphic or antiholomorphic bijection between domains.
\paragraph{Statement recorded in the supplied source.}
What geometric or conformal restrictions on $G_1,G_2$ follow from $H(G_1)\equiv H(G_2)$? Characterize elementary equivalence of these rings, in contrast with ring isomorphism, which the source identifies with conformal or anticonformal equivalence. Retain the annulus subcase $G_j=\{z:r_j<|z|<R_j\}$, $0<r_j<R_j<\infty$, singled out in the update. Results for the scalar-named complex algebra do not replace the ring-language question. \sref{2308005}{2}
\subsection{Short English statement}
Determine what elementary equivalence of holomorphic-function rings implies about their domains, including the annulus case and the distinction from complex-algebra equivalence.
\subsection{Sources}
\begin{description}
\item[\hypertarget{source:2308005:1}{S1}] ULAM.AI and the UnsolvedMath Contributors, UnsolvedMath: A Curated Collection of Open Mathematics Problems, record 2308005 (AMR-022-8005). CC BY 4.0. \url{https://huggingface.co/datasets/ulamai/UnsolvedMath}.
\item[\hypertarget{source:2308005:2}{S2}] Walter K. Hayman and Edward F. Lingham, Research Problems in Function Theory (2018), Problem 8.5, its complete update and relevant preceding definitions. \url{https://arxiv.org/abs/1809.07200}.
\end{description}
\label{end:2308005}
\end{document}
